\subsection{QUOTIENT GROUPS}

THEOREM 1. Let R be an equivalence relation on a group with operators G; if R is left (resp. right) compatible (§ 3, no. 3) with the group law on G and compatible with the action of Ω, the equivalence class of e is a stable subgroup H of G and the relation R is equivalent to \(x^{-1}y \in H\) (resp. \(yx^{-1} \in H\) ). Conversely, if H is a stable subgroup of G, the relation \(x^{-1}y \in H\) (resp. \(yx^{-1} \in H\) ) is an equivalence relation which is left (resp. right) compatible with the group law on G and compatible with the action of Ω and under which H is the equivalence class of e.

We restrict our attention to the case where the relation R is left compatible with the law on G (the case of a right compatible relation follows by replacing the law on G by the opposite law). The relation \(y \equiv x \pmod{R}\) is equivalent to \(x^{-1}y \equiv e \pmod{R}\) , for \(y \equiv x\) implies \(x^{-1}y \equiv x^{-1}x = e\) and conversely \(x^{-1}y \equiv e\) implies \(y = x(x^{-1}y) \equiv x\) . If H denotes the equivalence class of e, the relation R is then equivalent to \(x^{-1}y \in H\) . We show that H is a stable subgroup of G. For every operator \(\alpha\) , the relation \(x \equiv e\) implies \(x^{\alpha} \equiv e^{\alpha} = e\) , hence \(H^{\alpha} \subset H\) and H is stable under the action of \(\Omega\) . It suffices to establish (Proposition 1) that \(x \in H\) and \(y \in H\) imply \(x^{-1}y \in H\) , that is \(x \equiv e\) and \(y \equiv e\) imply \(x \equiv y\) , which is a consequence of the transitivity of R.

Conversely, let H be a stable subgroup of G; the relation \(x^{-1}y \in H\) is reflexive since \(x^{-1}x = e \in H\) ; it is symmetric since \(x^{-1}y \in H\) implies \(y^{-1}x = (x^{-1}y)^{-1} \in H\) ; it is transitive, for \(x^{-1}y \in H\) and \(y^{-1}z \in H\) imply \(x^{-1}z = (x^{-1}y)(y^{-1}z) \in H\) ; it is left compatible with the law of composition on G, for \(x^{-1}y = (zx)^{-1}(zy)\) for all \(z \in G\) ; finally, for every operator \(\alpha\) , the relation \(y \in xH\) implies \(y^{\alpha} \in x^{\alpha}H^{\alpha} \subset x^{\alpha}H\) and hence the equivalence relation \(x^{-1}y \in H\) is compatible with the action of \(\Omega\) on G.

Let G be a group and H a subgroup of G; the relation \(x^{-1}y \in H\) (resp. \(yx^{-1} \in H\) ) is also written in the equivalent form \(y \in xH\) (resp. \(y \in Hx\) ). Every subgroup H of G thus defines two equivalent relations on G, namely \(y \in xH\) and \(y \in Hx\) : the equivalent classes under these relations are respectively the sets xH, which are called left cosets of H (or modulo H), and the sets Hx, which are called right cosets of H (or modulo H). By saturating a subset A ⊂ G with respect to these relations (Set Theory, II, §6, no. 4), we obtain respectively the sets AH and HA. The mapping \(x \mapsto x^{-1}\) transforms left cosets modulo H into right cosets modulo H and conversely.

The cardinal of the set of left cosets (mod. H) is called the index of the subgroup H with respect to G and is denoted by (G:H); it is also equal to the cardinal of the set of right cosets.

If a subgroup K of G contains H, it is a union of left (or right) cosets of H. Since a left coset of K is obtained from K by left translation, the set of left cosets of H contained in a left coset of K has cardinal independent of the latter. Hence (Set Theory, III, § 5, no. 8, Proposition 9):

PROPOSITION 4. Let H and K be two subgroups of a group G such that H ⊂ K. Then

\[
(\mathbf {G}: \mathbf {H}) = (\mathbf {G}: \mathbf {K}) (\mathbf {K}: \mathbf {H}).\tag{2}
\]

COROLLARY. If G is a finite group of order g and H is a subgroup of G of order h, then

\[
h. (\mathbf {G}: \mathbf {H}) = g\tag{3}
\]

(in particular, the order and index of H are divisors of the order of G).

Theorem 1 allows us to determine the equivalence relations compatible with the laws on a group with operators G: if R is such a relation, it is both left and right compatible with the group law on G and with the action of \(\Omega\) . Hence, if H is the class of e (mod. R), H is a stable subgroup such that the relations \(y \in xH\) and \(y \in Hx\) are equivalent (since both are equivalent to R); hence \(xH = Hx\) for all \(x \in G\) . Conversely, if this is so, one or other of the equivalent relations \(y \in xH\) , \(y \in Hx\) is compatible with the group law, since it is both left and right compatible with this law (§ 3, no. 4) and is compatible with the action of \(\Omega\) . Since the equation \(xH = Hx\) is equivalent to \(xHx^{-1} = H\) , we make the following definition:

DEFINITION 5. Let G be a group with operators. A stable subgroup H of G is called a normal (or invariant) stable subgroup of G if \(xHx^{-1} = H\) for all \(x \in G\) .

If \(\Omega = \varnothing\) , a normal stable subgroup of G is called a normal (or invariant) subgroup of G. In a commutative group every subgroup is normal.

To verify that a stable subgroup H is normal, it suffices to show that \(xHx^{-1} \subset H\) for all \(x \in G\) ; for if so then \(x^{-1}Hx \subset H\) for all \(x \in G\) , that is \(H \subset xHx^{-1}\) , and hence \(H = xHx^{-1}\) .

Let H be a normal stable subgroup of G and R the equivalence relation \(y \in xH\) defined by H; on the quotient set G/R, the internal law, the quotient by R of the law of the group G, is associative; the class of e is the identity element of this quotient law; the classes of two inverse elements in G are inverses under the quotient law and the action of \(\Omega\) , the quotient by R of the action of \(\Omega\) on G, is distributive with respect to the internal law on G/R ( \(\S 3\) , no. 5). Hence, summarizing the results obtained:

THEOREM 2. Let G be a group with operators. For an equivalence relation R on G to be compatible with the group law and the action of \(\Omega\) , it is necessary and sufficient that it be of the form \(x^{-1}y \in H\) , where H is a normal stable subgroup of G (the relation \(x^{-1}y \in H\) being moreover equivalent to \(yx^{-1} \in H\) for such a subgroup). The law of composition on G/R the quotient of that on G and the action of \(\Omega\) on G/R the quotient of that of \(\Omega\) on G by such a relation R give G/R the structure of a group with operators, called the quotient structure, and the canonical mapping of the passage to the quotient is a homomorphism of groups with operators.

DEFINITION 6. The quotient of a group with operators G by the equivalence relation defined by a normal subgroup H of G, with the quotient structure, is called the quotient group with operators of G by H and is denoted by G/H. The canonical mapping G → G/H is called a canonical homomorphism

Let G be a group and H a normal subgroup of G. The quotient G/H, with its group structure, is called the quotient group of G by H. For a mapping from G/H to a group with operators to be a homomorphism of groups with operators, it is necessary and sufficient that its composition with the canonical mapping of G onto G/H be one: this justifies the name ``quotient group'' (Set Theory, IV, § 2, no. 6).

The equivalence relation defined by a normal stable subgroup of G is denoted by \(x \equiv y \pmod{H}\) or \(x \equiv y(H)\) .

PROPOSITION 5. Let \(f: G \to G'\) be a homomorphism of groups with operators and H and H' normal stable subgroups of G and G' respectively such that \(f(H) \subset H'\) . The mapping \(f\) is compatible with the equivalence relations defined by H and H'. Let \(\pi: G \to G/H\) and \(\pi': G' \to G'/H'\) be the canonical homomorphisms. The mapping \(\bar{f}: G/H \to G'/H'\) derived from \(f\) by passing to the quotients is a homomorphism.

If \(x \equiv y \pmod{H}\) , then \(x^{-1}y \in H\) , whence

\[
f (x) ^ {- 1} f (y) = f \left(x ^ {- 1}\right) f (y) = f \left(x ^ {- 1} y\right) \in f (\mathrm{H}) \subset \mathrm{H} ^ {\prime}
\]

and hence \(f(x) \equiv f(y) \pmod{H'}\) . The second assertion follows from the universal property of quotient laws (§ 1, no. 6).

Remarks. (1) If A is any subset of a group G and H is a normal subgroup of G, then AH = HA; this set is obtained by saturating A with respect to the relation \(x \equiv y \pmod{H}\) .

\begin{enumerate}
\def\labelenumi{(\arabic{enumi})}
\setcounter{enumi}{1}
\tightlist
\item
  If H is a normal subgroup of G of finite index, the quotient group G/H is a finite group of order (G:H).
\end{enumerate}

Note that if H is a normal subgroup of a group G and K is a normal subgroup of H, K is not necessarily a normal subgroup of G (I, § 5, Exercise 10).

Let G be a group with operators. The intersection of every family of normal stable subgroups of G is a normal stable subgroup. Hence, for every subset X of G, there exists a smallest normal stable subgroup containing X, called the normal stable subgroup generated by X.

In a group with operators G, the stable subgroups G and \(\{e\}\) are normal.

DEFINITION 7. A group with operators G is called simple if \(G \neq \{e\}\) and there exists no normal stable subgroup of G other than G and \(\{e\}\) .

